\documentclass[11pt]{article}
\usepackage[margin=1in]{geometry}
\usepackage{amsmath,amssymb,amsthm,mathtools}
\usepackage{hyperref}
\newtheorem{theorem}{Theorem}
\newtheorem{lemma}{Lemma}
\newtheorem{remark}{Remark}
\newcommand{\E}{\mathbb E}
\newcommand{\C}{\mathbb C}
\newcommand{\R}{\mathbb R}
\newcommand{\ip}[2]{\langle #1,#2\rangle}

\title{Degree-Four Vector-Coordinate Sum-of-Squares Cannot Exclude Seven MUBs in $\C^6$\\\large Verification note and exact certificate package}
\author{ProofNet24 P21 verification track}
\date{16 September 2026}

\begin{document}
\maketitle

\begin{abstract}
Randomstrasse101 Open Problem 23 asks whether degree-four Sum-of-Squares (SoS) can prove that seven mutually unbiased bases do not exist in $\C^6$.  For the two vector-coordinate encodings explicitly described on the source page, the answer is no.  This was proved by Sarkar in 2026.  We give a self-contained verification tailored to the source statement, resolve a normalization typo on the source page, and provide an exact rational/algebraic checker for the finite moment and localizer cores.  In the quartic equality formulation, the tracked 21 labeled cross-basis blocks are all simultaneously feasible under the degree-four Haar pseudoexpectation, so the number of exactly eliminated blocks is $0/21$.  This note does not address the unrestricted conjecture $\operatorname{MUB}(6)<7$ or stronger projector/noncommutative formulations.
\end{abstract}

\section{Source question and normalization}
Two orthonormal bases $\{v_i\}_{i=1}^d$ and $\{w_j\}_{j=1}^d$ are mutually unbiased when
\[
  |\ip{v_i}{w_j}|=\frac1{\sqrt d}\qquad\text{for all }i,j.
\]
Hence in dimension six the squared-overlap condition is
\[
  \boxed{|\ip{v_i}{w_j}|^2=\frac16.}
\]
The Randomstrasse101 page displays $1/\sqrt6$ inside one already-squared cross-overlap equation.  That display conflicts with the definition immediately above it and with the subsequent $2\times2$ LMI threshold $1/\sqrt6$.  Sarkar explicitly records the same typographical issue.  Throughout this note the quartic equality therefore uses $1/6$, while the amplitude/LMI threshold uses $1/\sqrt6$.

Fix $m=7$ labeled bases $B_k=(v_1^{(k)},\ldots,v_6^{(k)})$.  There are
\[
  \binom72=21
\]
labeled unordered cross-basis blocks and $6^2=36$ cross-overlap equations in each block, for $756$ quartic cross equations in total.

\section{What ``degree four'' means}
Work over the real polynomial ring in the real and imaginary parts of all vector coordinates.  A degree-four pseudoexpectation is a real linear functional
\[
 \widetilde\E:\R[x]_{\le4}\to\R
\]
such that $\widetilde\E[1]=1$ and
\[
 \widetilde\E[p^2]\ge0\qquad(\deg p\le2).
\]
For an equality $q=0$, the degree-four constraint is
\[
 \widetilde\E[hq]=0\quad\text{whenever }\deg(hq)\le4.
\]
Thus a quadratic same-basis orthonormality equation may be multiplied by a polynomial of degree at most two, whereas a quartic cross-unbiasedness equation may be multiplied only by a constant.

For a Hermitian matrix polynomial $M(x)\succeq0$ of degree at most two, degree-four localizing positivity requires
\[
 \widetilde\E[g(x)^*M(x)g(x)]\ge0
\]
for affine-linear vector polynomials $g$.  This is the convention used below for the source's $2\times2$ LMI encoding.

\section{Quartic equality formulation}
\begin{theorem}[Haar degree-four pseudoexpectation]
For every $d\ge2$ and every $m\ge1$, the direct vector-coordinate quartic equality formulation of $m$ mutually unbiased bases in $\C^d$ has a feasible degree-four pseudoexpectation.  Consequently it has no degree-four SoS refutation.
\end{theorem}

\begin{proof}
Let $U_1,\ldots,U_m$ be independent Haar-random unitary matrices in $U(d)$ and take $v_i^{(k)}=U_ke_i$.  For every real polynomial $p$ of degree at most four in the vector coordinates define
\[
 \widetilde\E[p]=\E[p(U_1,\ldots,U_m)].
\]
This is an actual expectation, so positivity on squares of degree at most two is automatic.

Within each basis, all orthonormality equations hold pointwise.  Therefore every allowed multiplier of a same-basis quadratic equality also has zero pseudoexpectation.

If $k\ne\ell$, then $x=v_i^{(k)}$ and $y=v_j^{(\ell)}$ are independent Haar unit vectors.  Unitary invariance gives
\[
 \E|\ip{x}{y}|^2=\frac1d.
\]
Hence
\[
 \widetilde\E\left[|\ip{v_i^{(k)}}{v_j^{(\ell)}}|^2-\frac1d\right]=0.
\]
The bracket is already quartic in the real vector coordinates, so at degree four it admits only constant multipliers.  Thus all degree-four constraints are satisfied.

If a degree-four SoS refutation existed, applying $\widetilde\E$ to the refutation identity would make every square term nonnegative and every equality multiple vanish, while the left side would equal $-1$, a contradiction.
\end{proof}

For $d=6,m=7$, this same pseudoexpectation satisfies all $756$ cross equations simultaneously.  Therefore the frozen metric is
\[
 \boxed{\text{exactly eliminated labeled cross-basis blocks}=0/21.}
\]
This is a nonexistence result for a degree-four refutation, not a failed search.

\section{Exact rational Gram certificate for one cross pair}
For one independent Haar pair $x,y\in\C^6$, identify each complex unit vector with a real vector on $S^{11}$.  Let $D=12$.  The exact real-sphere moments through degree four are
\[
 \E X_aX_b=\frac{\delta_{ab}}D,
\]
and
\[
 \E X_aX_bX_cX_e
 =\frac{\delta_{ab}\delta_{ce}+\delta_{ac}\delta_{be}+\delta_{ae}\delta_{bc}}{D(D+2)}.
\]
The checker accompanying this note builds the $325\times325$ moment Gram matrix for the degree-$\le2$ real monomial basis in $(x,y)$ and reconstructs it as a sum of $433$ positive rational weighted squares.  Entrywise equality is checked over \texttt{fractions.Fraction}.  It also verifies exactly
\[
 \E(\Re\ip{x}{y})^2=\frac1{12},\qquad
 \E(\Im\ip{x}{y})^2=\frac1{12},\qquad
 \E|\ip{x}{y}|^2=\frac16.
\]
The finite certificate is only a checker for the algebraic core; the theorem above follows from the full Haar expectation argument.

\section{The $2\times2$ Hermitian LMI formulation}
The source also proposes, for a cross pair $x,y$, the Hermitian constraint
\[
 M(z)=
 \begin{pmatrix}
 c&z\\ \bar z&c
 \end{pmatrix}\succeq0,
 \qquad z=\ip{x}{y},\qquad c=\frac1{\sqrt d}.
\]
Classically, together with orthonormality, the inequalities $|\ip{x}{y}|\le1/\sqrt d$ force equality of all squared overlaps because the $d$ squared overlaps of a unit vector with an orthonormal basis sum to one.

\begin{lemma}[Exact affine Haar compression]
Let $x,y$ be columns from two independent Haar bases and let $H$ be the complexified affine coordinate space generated by the constant and all entries and conjugate entries of all proposed Haar bases, with the $L^2$ inner product.  Define $T:H\to H$ by
\[
 \ip{f}{Tg}_{L^2}=\E[\bar f\,\ip{x}{y}\,g].
\]
Then
\[
 \|T\|_{\mathrm{op}}=\frac1d,
 \qquad \operatorname{rank}T=2d.
\]
\end{lemma}

\begin{proof}
Normalize affine coordinate modes by $\sqrt d$.  Write $z=\sum_{r=1}^d\bar x_r y_r$.  Haar second moments and independence give
\[
 T(\sqrt d\,\bar y_r)=\frac1d\sqrt d\,\bar x_r,
 \qquad
 T(\sqrt d\,x_r)=\frac1d\sqrt d\,y_r.
\]
All other affine modes vanish after projection.  Indeed, the constant channel has $\E z=0$; unmatched conjugation patterns are killed by independent column phases; coordinates belonging to other columns have a zero Kronecker-delta second moment; and coordinates from a third independent basis factor through a zero first moment.  The displayed $2d$ input modes and $2d$ output modes are orthonormal, so $T$ is exactly $1/d$ times a rank-$2d$ partial isometry.
\end{proof}

\begin{theorem}[Degree-four LMI feasibility]
The same Haar pseudoexpectation satisfies every degree-four localizing constraint for the source's $2\times2$ Hermitian vector-coordinate formulation.  Consequently this formulation also has no degree-four SoS refutation.
\end{theorem}

\begin{proof}
For affine-linear complex polynomials $p,q\in H$,
\[
\begin{aligned}
 \widetilde\E
 \begin{bmatrix}\bar p&\bar q\end{bmatrix}
 M(z)
 \begin{bmatrix}p\\q\end{bmatrix}
 &=c(\|p\|_2^2+\|q\|_2^2)+2\Re\ip{p}{Tq}_{L^2}\\
 &\ge c(\|p\|_2^2+\|q\|_2^2)-\frac2d\|p\|_2\|q\|_2\\
 &\ge\left(c-\frac1d\right)(\|p\|_2^2+\|q\|_2^2)\ge0,
\end{aligned}
\]
because $c=1/\sqrt d\ge1/d$.  Hence all degree-four localizers are positive.  Applying the pseudoexpectation to a putative mixed equality/LMI degree-four refutation again yields a contradiction.
\end{proof}

For $d=6,m=7$, the complex affine space has dimension
\[
 1+2md^2=505,
\]
so the localizer acts on a $1010$-dimensional space.  Since $T$ has rank $12$ and every nonzero singular value is $1/6$, a singular-vector basis reduces the localizer to twelve copies of
\[
 B=\begin{pmatrix}\sqrt6/6&1/6\\1/6&\sqrt6/6\end{pmatrix}
\]
and $986$ positive scalar directions.  Exactly,
\[
 \det B=\frac5{36}>0,
\]
and
\[
 B=
 \begin{pmatrix}1&0\\\sqrt6/6&1\end{pmatrix}
 \begin{pmatrix}\sqrt6/6&0\\0&5\sqrt6/36\end{pmatrix}
 \begin{pmatrix}1&\sqrt6/6\\0&1\end{pmatrix}.
\]
The checker verifies this factorization in exact arithmetic over $\mathbb Q(\sqrt6)$.

\section{Scope and literature status}
The conclusion is representation-sensitive.  In vector coordinates, cross-unbiasedness is quartic and degree four only tests its expectation.  In projector coordinates $P=vv^*$, the same relation becomes quadratic, so degree four can test substantially more.  Sarkar shows that a centered projector-coordinate degree-four Gram system recovers the elementary bound $m\le d+1$; for $d=6,m=7$ that bound is exactly tight and does not decide the unrestricted dimension-six MUB conjecture.

Accordingly, the result proved here is precisely:
\begin{quote}
For the two direct vector-coordinate degree-four formulations explicitly described in Randomstrasse101 Open Problem 23, there is no degree-four SoS refutation of seven proposed MUBs in $\C^6$.
\end{quote}
It does not prove that seven actual MUBs exist, and it does not rule out stronger projector, lifted, noncommutative, or higher-degree certificates.

The same source-matching theorem is already known from Shreyhaan Sarkar, \emph{Degree-Four Vector-Coordinate SoS Cannot Detect the MUB Upper Bound}, arXiv:2606.13903 (2026).  This document is therefore a verification/certificate package, not a claim of novelty.

\section{Reproducibility}
Run
\begin{verbatim}
python exact_checker.py
\end{verbatim}
from this folder.  The archived \texttt{exact\_checker.out} ends with
\begin{verbatim}
ALL_EXACT_CHECKS_PASS
\end{verbatim}
and records the exact metric $0/21$, the $325\times325$ rational moment-Gram identity, and the $\mathbb Q(\sqrt6)$ LMI block certificate.

\begin{thebibliography}{9}
\bibitem{Bandeira2025}
A. S. Bandeira, ``Mutually Unbiased Bases, ETFs, and Zauner's Conjecture (problems 22--24),'' Randomstrasse101, posted 14 December 2025, updated 18 December 2025.

\bibitem{Sarkar2026}
S. Sarkar, ``Degree-Four Vector-Coordinate SoS Cannot Detect the MUB Upper Bound,'' arXiv:2606.13903, 2026.

\bibitem{BDK2022}
A. S. Bandeira, N. Doppelbauer, and D. Kunisky, ``Dual bounds for the positive definite functions approach to mutually unbiased bases,'' Sampling Theory, Signal Processing, and Data Analysis 20 (2022), 18.

\bibitem{GP2024}
S. Gribling and S. Polak, ``Mutually unbiased bases: polynomial optimization and symmetry,'' Quantum 8 (2024), 1318.

\bibitem{MW2026}
D. McNulty and S. Weigert, ``Mutually Unbiased Bases in Composite Dimensions -- A Review,'' Quantum 10 (2026), 2051.
\end{thebibliography}

\end{document}
